\chapter{Erros comuns}
\label{cap:erros}

A tabela a seguir reúne as mensagens que mais aparecem para quem está
começando, com a causa e a correção de cada uma. Vale a pena consultá-la antes
de pedir ajuda: quase todo erro de aula está aqui.

\begin{longtable}{@{}p{4.6cm}p{4.4cm}p{4.6cm}@{}}
\caption{Erros frequentes, causas e correções}
\label{tab:erros} \\
\toprule
\textbf{Erro ou mensagem} & \textbf{Causa} & \textbf{Correção} \\
\midrule
\endfirsthead
\multicolumn{3}{@{}l}{\footnotesize\emph{Tabela~\ref{tab:erros} --- continuação}} \\
\toprule
\textbf{Erro ou mensagem} & \textbf{Causa} & \textbf{Correção} \\
\midrule
\endhead
\bottomrule
\endfoot
\footnotesize\comando{ERROR 1046: No database selected}
  & \footnotesize Faltou escolher o banco
  & \footnotesize \comando{USE liga\_SEUNOME;} ou clicar no banco na lateral \\
\addlinespace
\footnotesize\comando{ERROR 1044: Access denied to database}
  & \footnotesize Nome fora do padrão \comando{liga\_*}, ou tentativa de abrir o banco \comando{cap}
  & \footnotesize Use \comando{liga\_SEUNOME}: underscore, minúsculas, sem acento \\
\addlinespace
\footnotesize\comando{ERROR 1050: Table already exists}
  & \footnotesize O \comando{CREATE TABLE} foi executado duas vezes
  & \footnotesize \comando{DROP TABLE membro;} antes, ou rode o script do Apêndice~\ref{ap:script} inteiro \\
\addlinespace
\footnotesize\comando{ERROR 1062: Duplicate entry}
  & \footnotesize Violou \comando{PRIMARY KEY} ou \comando{UNIQUE}
  & \footnotesize O valor já existe; use outro ou faça \comando{UPDATE} \\
\addlinespace
\footnotesize\comando{ERROR 1452: foreign key constraint fails}
  & \footnotesize Chave estrangeira apontando para um \comando{id} inexistente
  & \footnotesize Insira primeiro na tabela-pai (\comando{curso}) \\
\addlinespace
\footnotesize\comando{ERROR 1452} ao apagar da tabela-pai
  & \footnotesize Há filhos referenciando aquela linha
  & \footnotesize Trate os filhos antes, ou defina \comando{ON DELETE} \\
\addlinespace
\footnotesize\comando{ERROR 1052: Column is ambiguous}
  & \footnotesize Duas tabelas com coluna de mesmo nome
  & \footnotesize Qualifique: \comando{m.nome}, \comando{c.nome} \\
\addlinespace
\footnotesize \comando{UPDATE} ou \comando{DELETE} afetou a tabela toda
  & \footnotesize \comando{WHERE} esquecido
  & \footnotesize Sempre testar com \comando{SELECT} antes (Seção~\ref{sec:where-ausente}) \\
\addlinespace
\footnotesize \comando{WHERE campo = NULL} não retorna nada
  & \footnotesize \comando{NULL} não é comparável com \comando{=}
  & \footnotesize \comando{WHERE campo IS NULL} \\
\addlinespace
\footnotesize Acentuação vira caractere estranho
  & \footnotesize Conjunto de caracteres errado (\comando{latin1})
  & \footnotesize Crie o banco com \comando{utf8mb4} \\
\addlinespace
\footnotesize Datas não filtram direito
  & \footnotesize Formato ou tipo errado
  & \footnotesize Use \comando{DATE}/\comando{DATETIME} e \comando{'AAAA-MM-DD'} \\
\addlinespace
\footnotesize Valores monetários erram centavos
  & \footnotesize Coluna declarada como \comando{FLOAT}
  & \footnotesize Use \comando{DECIMAL(10,2)} \\
\addlinespace
\footnotesize O comando não executa
  & \footnotesize Falta o ponto e vírgula final
  & \footnotesize Termine todo comando com \comando{;} \\
\end{longtable}
